% Solution d'éthanoate de sodium.

L'éthanoate de sodium $\ce{CH3COONa}$ est un produit dérivé de l'acide
éthanoïque $\ce{CH3COOH}$.

L'éthanoate de sodium est solide et se dissout dans l'eau en donnant une
solution constituée d'ions sodium et d'anions.

\begin{enumerate}
    \item Écrire l'équation-bilan de la réaction de dissolution dans l'eau
          (dissociation dans l'eau).
    \item Le pH de la solution obtenue est égal à 9. Que peut-on déduire
          concernant la solution~: est-elle acide~? neutre~? basique~? Justifier
          votre réponse.
    \item On ajoute à cette solution de l'acide chlorhydrique concentré
          (constitué d'ions $\ce{H3O^+}$ et d'ions $\ce{Cl^-}$).

          Comment évolue le pH du milieu réactionnel~? Écrire l'équation-bilan
          de la réaction qui a lieu.
\end{enumerate}
